% !TEX root = ../main.tex
\PlanPage{财务字段与清洗}{武桐宇}{23}

\textbf{本页结论：}财务清洗的目标不是生成更“漂亮”的数字，而是让评分输入能被追溯到合并口径、报告期、单位与处理规则；缺失或不一致字段应显式阻断计算。

现有数据字典约定财务金额以人民币百万元表示，并区分年度与季度字段\cite{wty-repo-dictionary}。进入评分前，必须保留原始值、清洗值和处理原因；累计值不能直接当作单季值，同一家公司不同报告期也不能在未说明时横向比较。当前演示字段为样例，不以此生成真实企业的财务结论。

\begin{table}[htbp]
\centering
\caption{财务输入字段的清洗与校验规则}
\label{tab:wty-financial-cleaning}
\small
\begin{tabularx}{\linewidth}{p{0.22\linewidth}p{0.23\linewidth}X}
\toprule
字段/问题 & 处理规则 & 输出限制 \\
\midrule
营业收入、毛利率 & 标注合并口径、报告期与币种；保留原始来源 & 无可比口径时不比较公司排名 \\
经营现金流 & 区分累计披露与单季值，必要时按相邻累计期差分 & 起始期或缺期不做差分推断 \\
资产负债率 & 检查分母、日期和异常范围，记录缺失与更正 & 不作为偿债能力评级 \\
负数与极端值 & 核对原始披露、单位和符号；不静默截断原值 & 异常无法解释时标记不可用 \\
日期错配 & 财务期、事件日和市场观察日分别存储 & 不将不同日期的字段拼成因果结论 \\
\bottomrule
\end{tabularx}
\end{table}

清洗日志应与数据包一起交付，包含字段名、旧值、新值、规则和操作者/脚本版本。这样复核者可重放转换过程，也可在原始数据被修订时定位受影响的评分。
